\subsubsection{A finite-order retraction suffices for finite wild image}

\begin{proposition}[Proved: finite-jet cone covers]
\label{D:finite-jet-cone-cover}
Let $Y\to S$ be a projective resolution of a normal excellent
surface germ over an algebraically closed field $k$ of
characteristic $p>0$.  Suppose it is an isomorphism off the
closed point and has sole reduced exceptional curve a smooth
proper connected curve $C$.  Put
\[
 I=\mathcal O_Y(-C),\qquad
 C_n=(C,\mathcal O_Y/I^n),\qquad
 N=N_{C/Y},\quad \deg N=-d<0,\quad g=g(C).
\]
Suppose for some integer $r\geq2$ that
\[
 C_r\longrightarrow C
 \quad\hbox{is a $k$-retraction of }C\hookrightarrow C_r,
 \qquad rd>2g-2.                                      \tag{JR}
\]
Choose a power $q=p^a$ with
\[
       qd>2g-2,\qquad q(rd-2g+2)>2g-2.                \tag{JQ}
\]
Then the completed germ admits a fixed finite local cone
cover of fraction-field degree $q$, whose radial valuation
extends that of $C$.  Its full exceptional wild inertia image
has order at most $q$.  The same bound holds for the
uncompleted germ and its strict henselization.
Only the $r$th infinitesimal neighborhood is assumed to
retract; no retraction of the entire formal neighborhood
is required.
\end{proposition}

\begin{proof}
Let $h:C'=C^{(-a)}\to C$ be the $k$-linear relative
Frobenius of degree $q$, and put $N'=h^*N$.
Use the retraction in (JR) to form the ordinary finite
scheme
\[
                      W_r=C_r\times_C C'.
\]
Its projection to $C_r$ is finite locally free and
radicial of degree $q$.  Let $J_r$ be its reduced
closed-curve ideal.  The powers of $J_r$ give
thickenings $W_n$, $1\leq n\leq r$, with
\[
 W_1=C',\qquad J_r^n/J_r^{n+1}=N'^{-n}
                    \quad(0\leq n<r).
\]
Here the equality at $n=0$ means the degree-zero
quotient $\mathcal O_{C'}$.  The retraction of $C_r$
base-changes to one of every $W_n$ onto $C'$.

Write $V=\operatorname{Tot}(N')$, with zero-section
ideal $J$, and $V_n=(C',\mathcal O_V/J^n)$.
We first construct a compatible isomorphism
$W_r\simeq V_r$ over $C'$, inducing the conormal identity.
On affine opens where $N'$ is trivial, the retraction
and a lift of a conormal generator identify $W_r$
with $\mathcal O_{C'}[s]/(s^r)$.
The first-order algebra is therefore split.
Suppose these coordinate identifications have been made
compatible modulo $s^n$, where $2\leq n<r$.
Their errors in degree $n$ are a cocycle in
\[
 \mathcal Hom(N'^{-1},N'^{-n})=N'^{\,1-n}.
 \tag{JC}
\]
Changing one coordinate by a term of degree $n$
adds a coboundary.  Since
\[
             \deg N'^{\,1-n}=qd(n-1)>2g-2,
\]
Serre duality kills this first cohomology group.
Induction for $n=2,\ldots,r-1$ gives $W_r\simeq V_r$;
when $r=2$ there are no such steps.
This is the finite version of the relative-coordinate
argument of Lemma~\ref{D:formal-retraction-cone-cover}.
There are no tangent-direction errors in (JC), since
all maps in this finite stage are maps over $C'$.

Composing the inverse of this isomorphism with
$W_r\to C_r\hookrightarrow Y$ gives
$f_r:V_r\to Y$.  Its closed-curve map is $h$, and its
conormal map is the identity
$h^*N^{-1}\simeq N'^{-1}$.  From this point onward we
use ordinary maps to the smooth surface $Y$, without
requiring them to be maps over a curve.

Suppose $n\geq r$ and $f_n:V_n\to Y$ has been constructed.
For the square-zero extension $V_n\subset V_{n+1}$
with ideal $N'^{-n}$, smoothness gives local lifts.
Their obstruction sheaf is
\[
                  h^*(T_Y|_C)\otimes N'^{-n}.
 \tag{JL}
\]
It has tangent and normal subquotients
$h^*T_C\otimes N'^{-n}$ and $N'^{\,1-n}$, as in
\eqref{D:cone-tangent}--\eqref{D:cone-normal}.
Their respective degrees are
\[
                 q(nd-2g+2),\qquad q(n-1)d.
 \tag{JD}
\]
For every $n\geq r$, the first is at least
$q(rd-2g+2)>2g-2$, and the second is at least
$qd>2g-2$.  Their $H^1$ groups vanish; the exact
tangent-normal sequence therefore kills $H^1$ of
(JL).  The local lifts glue to a global $f_{n+1}$.
It sends $I$ into $J$ and retains the conormal identity
already fixed at the finite stage.
The square-zero deformation and Serre-duality inputs
are the same precise primary statements used in
Lemma~\ref{D:finite-formal-cone-cover}.

The inverse limit of these compatible maps is a formal
map
\[
       \widehat{\operatorname{Tot}(N')}_{C'}
                     \longrightarrow\widehat Y_C.
\]
It is finite locally free of rank $q$: modulo a normal
parameter it is the finite flat map $h$, and the
parameter maps to a generator times a unit.  The
complete basis-lifting argument in
Lemma~\ref{D:finite-formal-cone-cover} applies directly.
Its algebraization proof, starting from this constructed
finite formal algebra, likewise applies without imposing
vanishing for the earlier values $2\leq n<r$.
It supplies an actual finite map over the completed
resolution, a normal local Stein cone algebra of
degree $q$, and its actual radial extension of the
exceptional DVR.

Finally use the inertia-index proof of
Proposition~\ref{D:finite-cone-wild-bound} for this one
fixed finite cone extension.  All original punctured
covers pull back to radially tame cone covers, so the
full wild image has order at most $q$.
That proposition's common-DVR comparison also gives
the assertion before completion and after strict
henselization.  The extension constructed after
order $r$ need not be purely inseparable, so this
argument proves finiteness of wild image, without
asserting that the image is zero.
\end{proof}

\begin{corollary}[Proved: a split first neighborhood]
\label{D:split-first-neighborhood-finite-wild}
Under the same surface hypotheses, suppose that
$C_2=(C,\mathcal O_Y/I^2)$ retracts onto $C$, or
equivalently that its square-zero class
$e_Y\in H^1(C,T_C\otimes N^{-1})$ is zero.
If $d>\max\{0,g-1\}$, the full exceptional wild
inertia image is finite.
One bound for its order is any power $q=p^a$ satisfying
\[
             qd>2g-2,\qquad q(2d-2g+2)>2g-2.
\]
\end{corollary}

\begin{proof}
The square-zero splitting equivalence follows from
the derivation-torsor calculation in
Lemma~\ref{D:finite-formal-cone-cover}.
Apply Proposition~\ref{D:finite-jet-cone-cover} with
$r=2$.  Its numerical condition is $2d>2g-2$,
which is exactly $d>g-1$ in addition to $d>0$.
\end{proof}
